\makeatletter
\appendix
% elsarticle makes \thesection expand to "Appendix~A", so \ref{} inside prose that
% already says "Appendix" renders "Appendix Appendix A". Number the appendices with
% bare letters and let the prose supply the word.
\gdef\thesection{\@Alph\c@section}
\makeatother

\section{Proofs}\label{app:proofs}

\subsection{Proof of Theorem~\ref{thm:population}}
Condition throughout on $C=c$. Expanding over labels gives
\[
\Delta T_c=\sum_y\mathbb E[H_X(y)\ind \{Y=y\}],\qquad
\Delta P_c=\sum_y\mathbb E[H_X(y)]\Pr(Y=y).
\]
Subtracting yields the covariance sum in Eq.~\eqref{eq:population-id}; the first
identity is the definition of $\Delta R_c$. Under the quadratic score,
\[
H_X(y)=2\Delta q_y(X)-\{\lVert q_1(\cdot\mid X)\rVert_2^2-
\lVert q_0(\cdot\mid X)\rVert_2^2\}.
\]
The bracketed term is independent of the label index. Its contribution summed over $y$
is $\operatorname{Cov}(B(X),\sum_y\ind \{Y=y\})=
\operatorname{Cov}(B(X),1)=0$. The remaining terms give Eq.~\eqref{eq:cov-id}. If
$H_X(y)=H_c(y)$ almost surely, every covariance is zero. \qed

\subsection{Proof of Proposition~\ref{prop:transported}}
Condition on $C=c$ and let $Y'\sim p(c)$ be independent of $X$. For the quadratic score
$S_{\rm B}(q,y)=2q_y-\lVert q\rVert^2$,
\begin{equation}
\mathbb E\bigl[S_{\rm B}(q_m,Y')\mid c\bigr]
=2\bigl\langle p(c),\bar q_m(c)\bigr\rangle-\mathbb E\bigl[\lVert q_m\rVert^2\mid c\bigr].
\end{equation}
Decompose the second moment as
$\mathbb E[\lVert q_m\rVert^2\mid c]=\lVert\bar q_m(c)\rVert^2
+\operatorname{tr}\operatorname{Var}(q_m\mid c)$, and complete the square using
$2\langle p,\bar q\rangle-\lVert\bar q\rVert^2=\lVert p\rVert^2-\lVert p-\bar q\rVert^2$:
\begin{equation}
\mathbb E\bigl[S_{\rm B}(q_m,Y')\mid c\bigr]
=\lVert p(c)\rVert^2-\lVert p(c)-\bar q_m(c)\rVert^2
-\operatorname{tr}\operatorname{Var}(q_m\mid c).
\end{equation}
Subtracting the $m=0$ expression from the $m=1$ expression, the $\lVert p(c)\rVert^2$
terms cancel and Eq.~\eqref{eq:transported-composition} follows, since
$\Delta P_c=\mathbb E[S_{\rm B}(q_1,Y')\mid c]-\mathbb E[S_{\rm B}(q_0,Y')\mid c]$ by
definition. \qed

\subsection{Proof of Theorem~\ref{thm:bregman}}
Condition on $C=c$. The proof of Theorem~\ref{thm:population} above establishes
$\Delta R_c=\sum_y\operatorname{Cov}(H_X(y),\ind\{Y=y\}\mid C=c)$ for \emph{any} score
contrast $H$, not only the quadratic one; only the identification of $H_x(y)$ with a
specific functional form is score-dependent. For the Bregman score of
the Bregman score $S_\psi(q,y)=\psi(q)+\langle\nabla\psi(q),e_y-q\rangle$,
\[
H_x(y)=\underbrace{\bigl[\psi(q_1(x))-\psi(q_0(x))\bigr]
-\bigl[\langle\nabla\psi(q_1(x)),q_1(x)\rangle-\langle\nabla\psi(q_0(x)),q_0(x)\rangle\bigr]}_{=:B(x)}
+\;\Delta g_y(x),
\]
writing $q_i(x)$ for $q_i(\cdot\mid x)$, where $B(x)$ collects every term not indexed by
$y$ and $\Delta g_y(x)=\partial_y\psi(q_1(x))-\partial_y\psi(q_0(x))$ is the $y$-indexed
remainder, exactly as in the proof of Theorem~\ref{thm:population} for the quadratic
case. Substituting,
\begin{align*}
\sum_y\operatorname{Cov}(H_X(y),\ind\{Y=y\}\mid c)
&=\operatorname{Cov}\Bigl(B(X),\sum_y\ind\{Y=y\}\ \Big|\ c\Bigr)\\
&\quad+\sum_y\operatorname{Cov}(\Delta g_y(X),\ind\{Y=y\}\mid c).
\end{align*}
The first term is $\operatorname{Cov}(B(X),1\mid c)=0$ since $\sum_y\ind\{Y=y\}=1$
almost surely. The second term is Eq.~\eqref{eq:bregman-cov-id}. \qed

\subsection{Proof of Theorem~\ref{thm:sample}}
Summing the kernel in Eq.~\eqref{eq:kernel} over unordered pairs, every observed term
$H_i(y_i)$ occurs in $n_c-1$ pairs with coefficient $1/2$, while the two crossed terms
together enumerate every ordered pair $H_i(y_j)$, $i\ne j$, with coefficient $1/2$.
Since $\binom{n_c}{2}=n_c(n_c-1)/2$,
\[
\widehat{\Delta R}_c
=\frac1{n_c}\sum_iH_i(y_i)
-\frac1{n_c(n_c-1)}\sum_{i\ne j}H_i(y_j)
=\widehat{\Delta T}_c-\widehat{\Delta P}_c.
\]
This is deterministic. When the cell's examples are i.i.d.\ draws from the cell
population, the average of the symmetric order-two kernel is the standard unbiased
U-statistic for its population expectation \cite{hoeffding1948class}; independence is
used, not merely exchangeability, because $\Delta P_c$ is defined through an independent
label coupling, so $\mathbb E[H_1(Y_2)]=\mathbb E[H_X(Y')]$ requires the two draws to be
independent. \qed

\subsection{Proof of Proposition~\ref{prop:coarsen}}
Fix $y$ and condition throughout on $\bar\phi=\bar c$. The law of total covariance
applied to the pair $(H_X(y),\ind \{Y=y\})$ with the finer conditioning variable $C$
nested inside $\bar c$ gives
\begin{align*}
\operatorname{Cov}\!\left(H_X(y),\ind \{Y=y\}\mid\bar\phi=\bar c\right)
&=\mathbb E\!\left[\operatorname{Cov}(H_X(y),\ind \{Y=y\}\mid C)
   \;\middle|\;\bar\phi=\bar c\right]\\
&\quad+\operatorname{Cov}\!\left(\mathbb E[H_X(y)\mid C],\,
   \Pr(Y=y\mid C)\;\middle|\;\bar\phi=\bar c\right),
\end{align*}
using $\mathbb E[\ind \{Y=y\}\mid C]=\Pr(Y=y\mid C)$. Summing over $y=1,\ldots,K$ and
applying the population identity of Theorem~\ref{thm:population} to the left side (at
$\bar\phi=\bar c$) and to the first term on the right side (at $C=c$, then averaged over
$c$ within $\bar c$) yields Eq.~\eqref{eq:coarsen}. \qed

\subsection{Proof of Proposition~\ref{prop:attenuate}}
Fix a cell $c$ with $n_c\ge2$ and $i\in c$. By definition of the label count,
$\sum_{y=1}^K m_{cy}H_i(y)=\sum_{j\in c}H_i(y_j)=H_i(y_i)+\sum_{j\ne i}H_i(y_j)$, and
Eq.~\eqref{eq:linear} gives $\sum_{j\ne i}H_i(y_j)=(n_c-1)\hat P_i$. Hence
\[
\tilde P_i=\frac1{n_c}\sum_{y}m_{cy}H_i(y)
=\frac{H_i(y_i)+(n_c-1)\hat P_i}{n_c}.
\]
Substituting $H_i(y_i)=\hat P_i+\hat R_i$ gives
$\tilde P_i=\hat P_i+n_c^{-1}\hat R_i$, and $\tilde R_i=H_i(y_i)-\tilde P_i
=\hat R_i-n_c^{-1}\hat R_i=(n_c-1)n_c^{-1}\hat R_i$. Averaging $\tilde R_i$ within cell
$c$ and reweighting by $n_c$ across cells, as in the dataset-level definition of
$\widehat{\Delta Z}$, gives
$\widetilde{\Delta R}=N_*^{-1}\sum_{c}n_c\cdot\frac{n_c-1}{n_c}\widehat{\Delta R}_c
=N_*^{-1}\sum_c(n_c-1)\widehat{\Delta R}_c$. \qed

\subsection{Proof of Proposition~\ref{prop:sensitivity}}
Fix $\phi=c$ throughout and write $a_y=\mathbb E[H_X(y)\mid\phi,Z]$,
$b_y=\Pr(Y=y\mid\phi,Z)$, both random through their dependence on $(\phi,Z)$. Let
$U=(a_y-\mathbb E[a_y\mid c])_{y=1}^K$ and $V=(b_y-\mathbb E[b_y\mid c])_{y=1}^K$ be the
corresponding mean-zero $\mathbb R^K$-valued random vectors conditional on $\phi=c$. The
map $(U,V)\mapsto\mathbb E[\langle U,V\rangle\mid c]=\sum_y\operatorname{Cov}(a_y,b_y\mid c)$
is an inner product on the space of mean-square-integrable $\mathbb R^K$-valued random
variables conditional on $\phi=c$, so by the Cauchy--Schwarz inequality applied to that
inner product,
\[
\Bigl|\sum_{y=1}^K\operatorname{Cov}(a_y,b_y\mid c)\Bigr|
\le\sqrt{\mathbb E[\lVert U\rVert^2\mid c]}\,\sqrt{\mathbb E[\lVert V\rVert^2\mid c]}
=\sqrt{\textstyle\sum_y\operatorname{Var}(a_y\mid c)}\,\sqrt{\textstyle\sum_y\operatorname{Var}(b_y\mid c)}.
\]
By the definition of $\rho$, the left side equals
$|\rho|\sqrt{\sum_y\operatorname{Var}(a_y\mid c)}\sqrt{\sum_y\operatorname{Var}(b_y\mid c)}$
exactly, so this step alone is not yet a bound; two further steps supply one, at two
different levels of conservatism. First, $a_y=\mathbb E[H_X(y)\mid\phi,Z]$ is itself a
conditional expectation of $H_X(y)$ given more information than $\phi=c$ alone, so the
law of total variance gives $\operatorname{Var}(H_X(y)\mid c)=\mathbb
E[\operatorname{Var}(H_X(y)\mid\phi,Z)\mid c]+\operatorname{Var}(a_y\mid c)\ge
\operatorname{Var}(a_y\mid c)$; likewise $b_y=\Pr(Y=y\mid\phi,Z)$ satisfies
$\operatorname{Var}(b_y\mid c)\le\operatorname{Var}(\ind\{Y=y\}\mid c)=p_y(c)(1-p_y(c))$ by
the same argument applied to $\ind\{Y=y\}$ in place of $H_X(y)$. Both bounds hold
termwise, so summing over $y$ and applying the (monotone) square root to each sum
separately gives the observable, per-cell bound
\[
\sqrt{\textstyle\sum_y\operatorname{Var}(a_y\mid c)}\sqrt{\textstyle\sum_y\operatorname{Var}(b_y\mid c)}
\le\sqrt{\textstyle\sum_y\operatorname{Var}(H_X(y)\mid c)}\sqrt{\textstyle\sum_yp_y(c)(1-p_y(c))},
\]
the second inequality of Eq.~\eqref{eq:sensitivity-bound}; this is a bound on a product of
two sums, and does not simplify to a sum of per-label square roots, since Cauchy--Schwarz
runs the other way for that rearrangement. Second, for
Corollary~\ref{cor:sensitivity-crude}, $a_y\in[-M,M]$ as a conditional expectation of a
$[-M,M]$-valued variable, so Popoviciu's inequality gives
$\operatorname{Var}(a_y\mid c)\le M^2$ for every $y$ and hence
$\sum_y\operatorname{Var}(a_y\mid c)\le KM^2$. For the $b_y$ the same argument would give
$\operatorname{Var}(b_y\mid c)\le1/4$ and a total of $K/4$, but that treats the $b_y$ as
$K$ unconstrained quantities in $[0,1]$ when in fact $\sum_yb_y=1$ almost surely. Using
that constraint, $\sum_yb_y^2\le\bigl(\sum_yb_y\bigr)^2=1$ pointwise, while
$\sum_y(\mathbb E b_y)^2\ge K^{-1}\bigl(\sum_y\mathbb E b_y\bigr)^2=1/K$ by
Cauchy--Schwarz, so
\[
\sum_{y=1}^K\operatorname{Var}(b_y\mid c)
=\mathbb E\Bigl[\sum_yb_y^2\Bigr]-\sum_y\bigl(\mathbb E b_y\bigr)^2\le1-\frac1K,
\]
which is Eq.~\eqref{eq:simplex-var} and is strictly smaller than $K/4$ for every $K>2$.
Therefore
\[
\Bigl|\sum_{y=1}^K\operatorname{Cov}(a_y,b_y\mid c)\Bigr|
\le|\rho|\sqrt{KM^2}\sqrt{1-1/K}=|\rho|\,M\sqrt{K-1}.
\]
By Eq.~\eqref{eq:sensitivity-exact}, the left side of each display is exactly
$|\Delta R_c(\phi)-\mathbb E[\Delta R_{(\phi,Z)}\mid\phi=c]|$, giving
Eq.~\eqref{eq:sensitivity-bound} and Eq.~\eqref{eq:sensitivity-crude}. \qed

\section{Influence function: estimand and derivation}\label{app:influence}

\paragraph{The estimand under singleton exclusion.} Cells with $n_c=1$ identify no
within-cell counterfactual and are excluded, so the dataset-level statistic targets the
\emph{identified-subpopulation} quantity: writing $\pi_c$ for the cell probabilities,
the target of $\widehat{\Delta R}$ is the $\pi$-weighted mean
$\theta=\sum_c\pi_c\,\Delta R_c/\sum_c\pi_c$ taken over cells that appear with at least
two examples. At finite $N$ the identified set is random; we treat the interval as
conditional on identification and report the identified fraction alongside every
estimate (99.0\% of relations on VG150, 100\% on CIFAR-100-LT), so the gap between the
conditional and full-population targets is bounded by the reported coverage.

\paragraph{Derivation of Eq.~\eqref{eq:influence}.} Fix a cell $c$ and let
$\theta_c=\Delta R_c$ with $g_c(z)=\mathbb E[A(z,Z')\mid Z'\in c]$ the kernel
projection. The classical H\'ajek projection of the order-two U-statistic
$\widehat{\Delta R}_c$ is
\[
\widehat{\Delta R}_c=\theta_c+\frac{2}{n_c}\sum_{i\in c}\bigl(g_c(Z_i)-\theta_c\bigr)
+O_p(n_c^{-1}).
\]
The dataset statistic weights cells by their example counts,
$\widehat{\Delta R}=N_*^{-1}\sum_c n_c\widehat{\Delta R}_c$, and cell membership is
itself random (multinomial over cells), so an example $i$ contributes both through the
within-cell projection and through the between-cell composition term
$\theta_{C_i}-\theta$. Combining the two sources,
\[
\widehat{\Delta R}-\theta
=\frac{1}{N_*}\sum_{i}\Bigl[2\bigl(g_{C_i}(Z_i)-\theta_{C_i}\bigr)
+\bigl(\theta_{C_i}-\theta\bigr)\Bigr]+o_p(N_*^{-1/2}),
\]
where the dataset-level remainder requires one condition beyond the per-cell expansion,
because $C$ per-cell remainders do not aggregate for free. Writing
$C=\#\{c:n_c\ge2\}$ for the number of eligible cells, the weighted sum of the per-cell
degenerate terms is $N_*^{-1}\sum_cn_c\cdot O_p(n_c^{-1})=O_p(C/N_*)$, which is
$o_p(N_*^{-1/2})$ provided
\begin{equation}
C=o\bigl(\sqrt{N_*}\bigr).
\label{eq:cellcount}
\end{equation}
This is what makes the second-order rate $O_p(n_c^{-1})$ rather than $o_p(n_c^{-1/2})$
worth stating: the sharper per-cell rate is what allows a condition on the cell
\emph{count} alone, with no further requirement on individual cell sizes. On VG150,
$C=6{,}346$ against $N_*=227{,}337$, so $C/\sqrt{N_*}\approx13.3$ --- the condition is a
statement about the limit, and at these finite sample sizes it is the coverage simulation
of Section~\ref{sec:simulation}, not Eq.~\eqref{eq:cellcount}, that supports the reported
intervals.
Estimating $g_{C_i}(Z_i)$ by the pair-kernel mean $\bar A_i$, $\theta_c$ by
$\widehat{\Delta R}_c$, and $\theta$ by $\widehat{\Delta R}$ gives the plug-in
influence contribution
$\widehat\psi_i=2(\bar A_i-\widehat{\Delta R}_{c})+(\widehat{\Delta R}_{c}
-\widehat{\Delta R})=2\bar A_i-\widehat{\Delta R}_{c}-\widehat{\Delta R}$, which is
Eq.~\eqref{eq:influence} and has exactly zero sample mean by construction. Its empirical
95\% coverage in a VG-like regime (image-clustered examples, heavy-tailed cell sizes with
a large mass of near-minimum cells) is validated in Section~\ref{sec:simulation}.

\paragraph{Proof of Theorem~\ref{thm:clt}.} The linearization above gives
$\widehat{\Delta R}-\theta=N_*^{-1}\sum_i\psi_i+o_p(N_*^{-1/2})$, under
Eq.~\eqref{eq:cellcount}, for the population
influence contribution $\psi_i=2(g_{C_i}(Z_i)-\theta_{C_i})+(\theta_{C_i}-\theta)$, which
has zero mean by construction and satisfies $|\psi_i|\le C(M)$ for an absolute constant
$C(M)$ depending only on the score bound $M$ in condition (i), since $g_{C_i}$,
$\theta_{C_i}$, and $\theta$ are themselves averages of the bounded kernel
$A_{ij}\in[-2M,2M]$. Because clusters are mutually independent (condition (ii)),
$T_g=\sum_{i\in g}\psi_i$ for $g=1,\dots,G$ are independent, mean-zero random variables
with $|T_g|\le|g|\,C(M)$. Condition (ii)'s finite third moment for cluster sizes and
$\max_g|g|=o(\sqrt G)$ give
$\sum_gE|T_g|^3\le C(M)^3\sum_g\mathbb E|g|^3=O(G)$ while
$\bigl(\sum_g\operatorname{Var}(T_g)\bigr)^{3/2}=\Theta(G^{3/2})$ under condition (iv), so
Lyapunov's ratio $\sum_gE|T_g|^3/(\sum_g\operatorname{Var}(T_g))^{3/2}\to0$; the
Lyapunov condition implies the Lindeberg condition, and the Lindeberg--Feller central
limit theorem for triangular arrays of independent summands
\citep{serfling1980approximation,vandervaart1998asymptotic} gives
$\bigl(\sum_gT_g\bigr)/\sqrt{\sum_g\operatorname{Var}(T_g)}\xrightarrow{d}\mathcal N(0,1)$.
Condition (iii) makes $N_*/\sum_g\operatorname{Var}(T_g)\to1/\sigma^2$ in probability by
definition of $\sigma^2$, and the $o_p(N_*^{-1/2})$ remainder is asymptotically
negligible after multiplying by $\sqrt{N_*}$, so
$\sqrt{N_*}(\widehat{\Delta R}-\theta)\xrightarrow{d}\mathcal N(0,\sigma^2)$ by Slutsky's
theorem. For the variance estimator, $\widehat{\Delta R}\to_p\theta$ by the weak law of
large numbers across the diverging number of identified examples, and
$\widehat{\Delta R}_c\to_p\Delta R_c$ for each eligible cell by the weak law of large
numbers applied within that cell -- this is where condition (iv) is used, and it is not
implied by anything aggregate: (iv) guarantees $n_c\to\infty$ almost surely for every cell
in the fixed, finite support of $\phi$, which is what a within-cell WLLN requires, whereas
a condition on the total identified count bounds no individual cell's size. So $\widehat\psi_i\to_p\psi_i$ pointwise, and
$\widehat\sigma^2$ -- a sum of squares of cluster-aggregated $\widehat\psi_i$, normalized
by $N_*$ and $G$ -- is a continuous functional of these consistent plug-ins, hence
$\widehat\sigma^2\to_p\sigma^2$ by the continuous mapping theorem. \qed

\paragraph{Efficient computation.} The implementation does not instantiate all pairs. In addition to Eq.~\eqref{eq:linear},
let $G_c(y)=\sum_{j\in c}H_j(y)$ and $O_c=\sum_{j\in c}H_j(y_j)$. The pair-kernel mean
for observation $i$ is
\[
\bar A_i=\frac12\left[
H_i(y_i)+\frac{O_c-H_i(y_i)}{n_c-1}
-P_i-\frac{G_c(y_i)-H_i(y_i)}{n_c-1}
\right].
\]
Therefore both the point estimate and H\'ajek projection require only label counts,
column sums, and matrix--vector products inside each cell. For image clusters $g$, the
implemented variance is
\[
\widehat{\operatorname{Var}}(\widehat{\Delta R})
=\frac{G}{G-1}\frac1{N_*^2}\sum_{g=1}^G
\left(\sum_{i\in g}\widehat\psi_i\right)^2,
\]
where $G$ is the number of images and $\widehat\psi_i$ is
Eq.~\eqref{eq:influence}. We use a normal critical value; with 27,955 image clusters,
small-$G$ corrections are immaterial.

\section{Randomization algorithm}\label{app:random}

For each cell, sort examples once and index them $0,\ldots,n_c-1$. Randomization $b$
draws $o_c^{(b)}$ uniformly from this index set and maps label position $r$ to
$(r+o_c^{(b)})\bmod n_c$. Independent offsets across cells define an element of the
product of cyclic groups. The cell label histogram is unchanged, so
$\sum_y m_{cy}H_i(y)$ in Eq.~\eqref{eq:linear} is precomputed; only the assigned label
lookup changes. One randomization is $O(N)$ after the $O(NK)$ setup. The plus-one Monte
Carlo correction includes the observed assignment and prevents zero $p$-values.

\section{Implementation and reproducibility}\label{app:implementation}

The NumPy implementation is in \texttt{wager/antisymmetric.py}. The primary routine
\texttt{decompose\_gain} validates probability matrices, forms score contrasts, excludes
singleton cells, computes all three components, verifies the identity numerically, and
returns per-instance transported/covariance values plus intervals. The Visual Genome
driver is \texttt{experiments/\allowbreak run\_sgg\_wager.py}; controlled validation and sensitivity
checks are in \texttt{experiments/\allowbreak antisymmetric\_simulation.py} and
\texttt{experiments/\allowbreak antisymmetric\_ablations.py}. All random seeds are fixed. Unit tests
cover the decomposition identity, quadratic covariance identity, cell-constant null,
singleton reporting, cluster inference, randomization power, and log-score path, plus
direct numerical checks of Propositions~\ref{prop:attenuate} and~\ref{prop:coarsen}
(the latter against an explicit discrete population with exact expectations, independent
of the estimator's own code path), Corollary~\ref{cor:bregman}'s log-score covariance
identity, Corollary~\ref{cor:dm}'s equivalence to the (unclustered) Diebold--Mariano
statistic at a trivial grouping variable, and Proposition~\ref{prop:sensitivity}'s bound
ordering (exact bias $\le$ tight, per-cell bound $\le$ crude, data-free bound) on a second
explicit discrete population with a known omitted confounder. The covariance-identity
cross-check quoted in
Appendix~\ref{app:lt} is reproduced by
\texttt{experiments/\allowbreak cifar\_covariance\_check.py}; the calibration-matched
protocol and the prior-matching experiment are
\texttt{experiments/\allowbreak cifar\_recalibration\_control.py} and
\texttt{experiments/\allowbreak vg\_prior\_consequence.py}. Table~\ref{tab:sensitivity-numeric}'s
free-bound rows are reproduced by \texttt{experiments/\allowbreak
sensitivity\_bound\_check.py} and its sharp rows, including the robustness value
$\rho^\dagger$ of Eq.~\eqref{eq:robustness-value}, by
\texttt{experiments/\allowbreak sensitivity\_analysis.py}; the evaluation of
Corollary~\ref{cor:sensitivity-crude} is
\texttt{experiments/\allowbreak sensitivity\_bound\_check.py}, which recomputes the crude
bound directly from the committed \texttt{antisymmetric\_ablations.json} rather than
re-running the decomposition; \texttt{experiments/\allowbreak sensitivity\_analysis.py}
computes Proposition~\ref{prop:sensitivity}'s tighter, per-cell form from per-instance
arrays (Section~\ref{app:compute}).

\paragraph{Full-data VG predictors (Appendices~\ref{app:vgmain}--\ref{sec:ablations}).}
FREQ uses add-one smoothing over the $50$ predicates within each training class pair,
backing off to the subject-marginal and then the global predicate histogram for pairs
unseen in training. The MLPs consume fixed random class embeddings
($32$ dimensions per class, seed $0$, concatenated for subject and object) on
standardized features and train with AdamW (learning rate $10^{-3}$, weight decay
$10^{-5}$, batch size $4096$): MLP-CLASS with hidden sizes $(256,128)$ for $12$ epochs,
MLP-SPATIAL with $(256,128)$ for $15$ epochs, and MLP-SPATIAL+ with $(512,256,128)$ for
$20$ epochs -- the ``+'' variant differs in both width and schedule, as noted in
Appendix~\ref{app:vgmain}.

\paragraph{Data availability.} The repository ships the estimator, every driver script,
and the cached results files that every printed number traces to
(\texttt{results/*.json}). The cached per-relation probability arrays that the drivers
consume (\texttt{sgg\_dists.npz}, \texttt{vg\_visual\_models.npz},
\texttt{cifar\_lt\_results.npz}) are regenerable from the committed training scripts
with fixed seeds and are archived alongside the code release, so exact reproduction
does not require retraining.

\subsection{Compute environment for Appendices~\ref{app:vgvisual}--\ref{app:lt}}\label{app:compute}

Model training and feature extraction for both new experiments ran on a single NVIDIA
T4 (15\,GB) GPU. The WAGER decomposition itself is pure NumPy and requires no
accelerator: it consumes the resulting cached probability arrays, as in
Appendix~\ref{app:vgmain}. All Python code runs under a pinned environment (Python 3.13).

\paragraph{Matched-subsample models (Appendix~\ref{app:vgvisual}).} The three compared
predictors share a fixed, seeded ($\text{seed}=0$) subsample of $100{,}000$ of the
$532{,}987$ training relations, drawn without replacement uniformly at the relation
level (not the image level); the full $229{,}605$-relation test split is always used for
evaluation. All three use the identical head and optimizer -- AdamW, learning rate
$10^{-3}$, weight decay $10^{-5}$, 15 epochs, batch size 4096, hidden sizes
$(256,128)$, on standardized features -- so they differ only in their input features.

Subject/object crops are taken from the annotated boxes and CLIP-preprocessed (bicubic
resize to $224\times224$ with CLIP's published normalization constants); a crop with no
valid overlap with the image falls back to a blank crop rather than being dropped, so
every relation contributes a feature row. CLIP runs in half precision with a batch size
of $512$ crops and is never fine-tuned. Two implementation details are worth recording.
First, the pretrained OpenAI weights were trained with the QuickGELU activation, so the
encoder must be instantiated as \texttt{ViT-B-32-quickgelu}; loading the same weights
under the plain \texttt{ViT-B-32} configuration silently substitutes a standard GELU and
degrades every embedding. Second, because an annotated object box participates in
several relations, crops are deduplicated by (image, box) before encoding and gathered
afterwards: the $659{,}210$ crops implied by the audited relations collapse to
$422{,}143$ unique ones, removing $36.0\%$ of the encoder work without changing any
feature. Only the images referenced by those relations are retrieved, individually rather
than through the two bulk archives; of the $71{,}990$ required, three were unavailable
upstream and fall back to a blank crop.

\paragraph{ResNet-32 triple (Appendix~\ref{app:lt}).} CIFAR-100-LT uses the standard
exponential-profile construction of \citet{cui2019classbalanced} at imbalance ratio
$100$: class $c$ (zero-indexed, $0,\ldots,99$) keeps
$n_c=\lfloor 500\cdot\mu^{c}\rfloor$ training images, where
$\mu=100^{-1/99}$, so class $0$ keeps all $500$ images and class $99$ keeps $5$; the
balanced $10{,}000$-image test split is untouched. All three models are a from-scratch,
randomly initialized CIFAR ResNet-32 (three stages of five basic residual blocks,
$16/32/64$ channels) trained for $80$ epochs with SGD (momentum $0.9$, Nesterov,
weight decay $2\times10^{-4}$), batch size $128$, standard crop-and-flip augmentation,
and a cosine learning-rate schedule (peak $0.1$, five-epoch linear warmup). The only
difference between the runs is the per-example loss weight. CE uses uniform weights
throughout. CB applies the effective-number class-balanced weight
$w_c\propto(1-\beta)/(1-\beta^{n_c})$ with $\beta=0.9999$, renormalized to average to
one across classes, from initialization. DRW trains with uniform weights until epoch
$60$ and switches to exactly the CB weights for the final twenty epochs
\citep{cao2019learning}; no other hyperparameter changes at the switch. The robustness
matrix of Table~\ref{tab:cifar-robust} repeats this recipe verbatim at seeds $1$ and $2$
(ratio $100$) and at ratios $10$ and $50$ (seed $0$), re-drawing the long-tailed subset
with the configuration's own seed and initializing all three arms of a configuration
identically (\texttt{experiments/colab\_cifar\_multiseed.py}). The post-hoc LA arm is not
trained: it rescales the baseline's test probabilities by the inverse training class
prior, $q_{\mathrm{LA}}(y\mid x)\propto q_{\mathrm{CE}}(y\mid x)/\pi_{\mathrm{train}}(y)$,
renormalized.

\section{Sensitivity bounds on the reported estimates}\label{app:sensitivity}

\paragraph{What the bound says about this paper's own estimates.}
Table~\ref{tab:sensitivity-numeric} reports both forms of the bound on the Visual Genome
pair of Appendix~\ref{sec:results}, together with the coarsening from the class-pair to
the subject-only cell, which instantiates Eq.~\eqref{eq:sensitivity-exact} exactly with
$Z=$ object class. The free bound is directionally correct and, at $K=50$, still
$1{,}887$ times the actual coarsening bias, so any $\bar\rho$ an analyst would plausibly
declare certifies a limit far above what is observed. That looseness is a property of a
worst-case inequality applied to label probabilities that concentrate away from the
extremes at which it binds, not a defect of this dataset --- but it also means the free
bound cannot support a claim about whether any particular reported $\widehat{\Delta R}$
would survive confounding.

The sharp per-cell form of Eq.~\eqref{eq:sensitivity-bound}, which replaces the
worst-case sums by the cell's own $\sum_y\operatorname{Var}(H_X(y)\mid c)$ and
$\sum_yp_y(c)(1-p_y(c))$, can support such a claim, and it is less reassuring than the
free bound. Computed by \texttt{experiments/sensitivity\_analysis.py} from per-instance
arrays for the same pair, it gives $B=0.23284$ at $\bar\rho=1$, so the
\emph{robustness value} --- the smallest aggregate correlation at which the bound reaches
the estimate --- is
\begin{equation}
\rho^\dagger=\frac{\bigl|\widehat{\Delta R}\bigr|}{B}
=\frac{0.01427}{0.23284}=0.06127 .
\label{eq:robustness-value}
\end{equation}
An unrecorded covariate whose aggregate correlation with the within-cell gain and the
within-cell label distribution reaches $0.061$ is enough to account for the whole of that
$\widehat{\Delta R}$. At the $\bar\rho=0.1$ that this section suggests as a plausible
elicited value, the sharp bound is $0.02328$, which \emph{exceeds} the estimate. We
therefore do not claim that a confounder would have to be implausibly strong to explain
these Visual Genome gains away; on this evidence a modest one would suffice, and the
reader should treat Appendix~\ref{sec:results}'s covariance gains as conditional on the
declared $\phi$ being close to complete. The flagship audit of
Section~\ref{sec:sggaudit} is not exposed in the same way, because its reported
$\widehat{\Delta R}$ is indistinguishable from zero and a sensitivity bound on a null
result changes nothing about it.

\begin{table}[t]
\centering
\caption{Both forms of the Proposition~\ref{prop:sensitivity} bound on the
MLP-SPATIAL vs.\ MLP-CLASS pair, and the coarsening already reported in
Table~\ref{tab:sensitivity} ($K=50$, $M=2$ for the quadratic score). The empirical bias
is the law-of-total-expectation aggregate of Eq.~\eqref{eq:sensitivity-exact}'s covariance
term across all subject-only cells.
\texttt{experiments/sensitivity\_bound\_check.py} reproduces the free-bound rows from the
committed ablation results and \texttt{experiments/sensitivity\_analysis.py} the sharp
rows from per-instance arrays.}
\label{tab:sensitivity-numeric}
\begin{tabular}{lr}
\toprule
Quantity & Value \\
\midrule
$\widehat{\Delta R}$, $\phi=$ subject-only & $0.02181$ \\
$\widehat{\Delta R}$, $\phi=$ class-pair & $0.01439$ \\
Empirical coarsening bias & $0.00742$ \\
\midrule
Free bound, $M\sqrt{K-1}$, at $\bar\rho=1$ & $14.00$ \\
Minimum $\bar\rho$ for the free bound to reach the bias & $0.00053$ \\
\midrule
Sharp bound $B$ at $\bar\rho=1$ & $0.23284$ \\
Sharp bound at $\bar\rho=0.1$ & $0.02328$ \\
Robustness value $\rho^\dagger$ & $0.06127$ \\
\bottomrule
\end{tabular}
\end{table}

\section{Conditions for the limit theorem}\label{app:clt}

Two things about condition (iv) deserve saying plainly rather than in a footnote. It is
what licenses treating the per-cell plug-ins $\widehat{\Delta R}_c$ as consistent for
$\Delta R_c$ in the variance argument, and a cell's size is not controlled by the
aggregate identified count, so this needs its own divergence requirement. But it also
forces every cell to be eventually identified, hence $N_*/N\to1$: an earlier version of
this theorem additionally assumed an identified fraction converging to some
$\kappa\in(0,1]$, which condition (iv) makes vacuous, and under which the
identified-subpopulation estimand $\theta$ would in any case coincide with the
unconditional target. We have therefore dropped that assumption rather than state two
conditions that cannot both bind. The price is that the theorem does not describe the
regime the paper's own applications sit in, where $\theta$ is a random, $N$-dependent
target because coverage is short of one ($99.0\%$ on VG150, $98.9\%$ in
Section~\ref{sec:sggaudit}). What the theorem licenses is the interval's form and its
variance estimator; what licenses reading the reported intervals at finite $N$ is the
coverage simulation of Section~\ref{sec:simulation}, and Section~\ref{sec:discussion}
states the gap between the two rather than eliding it.

The proof (Appendix~\ref{app:proofs}) combines the exact linearization of
Appendix~\ref{app:influence} with a Lindeberg--Feller central limit theorem for the
independent-across-cluster sum of mean-zero, uniformly bounded terms
\citep{serfling1980approximation,vandervaart1998asymptotic}; boundedness (i) makes the
Lindeberg condition automatic given (ii), and consistency of $\widehat\sigma^2$ follows
from the weak law of large numbers applied to each plug-in $\widehat{\Delta R}_c,
\widehat{\Delta R}$ together with continuity of the variance functional. One step needs
more care than the appendix's per-cell statement supplies: a per-cell H\'ajek remainder of
order $o_p(n_c^{-1/2})$ does not aggregate to $o_p(N_*^{-1/2})$ at the dataset level when
the number of cells grows, since summing $C$ such remainders admits a $\sqrt C$ inflation.
The repair does not need a new condition on cell sizes, only that the cell count be
negligible relative to the identified sample, $\#\{c:n_c\ge2\}=o(N_*)$: the second-order
degenerate term of a two-sample U-statistic is $O_p(n_c^{-1})$ rather than
$o_p(n_c^{-1/2})$, and $\sum_c n_c\cdot O_p(n_c^{-1})=O_p(C)=o_p(N_*)$ gives the
dataset-level rate directly. We state that condition explicitly; on VG150 it holds
comfortably ($C=6{,}346$ against $N_*=227{,}337$). Condition (ii) is what rules out a single cluster dominating the
sum -- an image contributing a positive fraction of all relations, for instance -- and
is a standard requirement for cluster-robust inference generally, not specific to this
estimator. This theorem licenses reading the image-clustered intervals reported
throughout Section~\ref{sec:exp} as asymptotically valid in the idealized regime where
every eligible cell's size grows without bound. That regime is not attained in-sample:
VG150's eligible cells are heavy-tailed, and $2{,}286$ of the $6{,}346$ eligible
class-pair cells hold fewer than five examples (Table~\ref{tab:sensitivity}). The
proportion is smaller than earlier drafts of this paper asserted --- those cells carry
$2.7\%$ of identified relations, and the median eligible cell holds appreciably more than
the minimum --- but a third of the cells still sit where condition (iv)'s divergence is
plainly hypothetical. The reported $94.0\%$ empirical coverage against a nominal $95\%$
(Section~\ref{sec:simulation}), not this theorem directly, is therefore the evidence that
the interval remains a good approximation well short of the idealization.

\section{Choice of the grouping variable}\label{app:partition}

\paragraph{Fineness is not a total order, and the guidance needs that qualification.}
The advice to prefer the finest defensible $\phi$ is sound against \emph{coarsening},
which is what Proposition~\ref{prop:coarsen} governs. It gives no guidance at all between
two partitions of the same granularity, neither of which coarsens the other --- and there
the estimand can differ without limit. Four cases suffice. Take two labels, $q_0$
uninformative at $(\tfrac12,\tfrac12)$, and a new model that is confidently right on all
four: $q_1=(0.9,0.1)$ on the two cases labelled $1$ and $(0.1,0.9)$ on the two labelled
$2$. The total gain is $\widehat{\Delta T}=+0.480$ under every partition, and every
partition below is fully identified at $100\%$ coverage:
\begin{center}
\begin{tabular}{@{}lccr@{}}
\toprule
declared $\phi$ & cells & $\widehat{\Delta P}$ & $\widehat{\Delta R}$ \\
\midrule
$\{1,3\},\{2,4\}$ (labels differ within each cell) & $2\times2$ & $-1.120$ & $+1.600$ \\
$\{1,2,3,4\}$ (one cell)                           & $1\times4$ & $-0.587$ & $+1.067$ \\
$\{1,2\},\{3,4\}$ (labels agree within each cell)  & $2\times2$ & $+0.480$ & $\phantom{+}0.000$ \\
\bottomrule
\end{tabular}
\end{center}
The first and third partitions are equally fine, both fully identified, and they disagree
about whether this model pair's improvement is entirely case-level or entirely
transported. The third is the degenerate case the paper already warns about --- grouping
by something that determines the label leaves each cell label-homogeneous and drives
$\widehat{\Delta R}$ to zero algebraically --- but the warning was previously attached to
the choice $\phi=Y$, and the point here is more general and less comfortable: any $\phi$
correlated with the label inside a cell partially degenerates the estimand, ``finest
defensible'' cannot detect it, and Proposition~\ref{prop:coarsen} is silent because
neither partition is a coarsening of the other. What the analyst actually needs to
declare, and what we recommend reporting alongside $\phi$, is the within-cell label
heterogeneity that remains: a partition whose cells are nearly label-homogeneous cannot
identify much of anything, and the coverage statistic does not reveal it. This is a
limitation of the construction rather than of any application of it, and we do not have a
procedure that resolves it.

\paragraph{A recalibration-invariant alternative, and what it would cost.}
Because $\Delta R$ is a covariance between forecast values, its sensitivity to
recalibration is structural rather than incidental, and a reader may reasonably ask why we
do not use a rank-based within-cell contrast instead --- a within-cell paired AUC
difference, say, for which \citet{delong1988comparing} already supply the
correlated-U-statistic variance. Such a quantity would be invariant to any monotone
rescaling of either model and would dissolve the confound this paper spends a protocol on.
We think the trade is worth stating and is not obviously in our favour on every axis.
What a rank-based contrast gives up is the exact identity: it is not a component of the
score difference, so it cannot be added to anything to recover $\Delta T$, and a benchmark
reporting it learns how the models rank cases but not how the reported score gain
decomposes. The construction here buys exact additivity with respect to a declared proper
score, and pays for it with scale dependence. A rank-based diagnostic is the natural
companion rather than a competitor, and reporting both --- one exact and scale-dependent,
one invariant and non-additive --- is what we would now recommend to a benchmark that
could choose.

% --- statements relocated from the method section for length ---
\begin{proposition}[Sensitivity bound for an unrecorded confounder]\label{prop:sensitivity}
Suppose $H_x(y)\in[-M,M]$ for all $x,y$ (WAGER already requires a bounded score). Fix a
cell $c$ of $\phi$, write $a_y=\mathbb E[H_X(y)\mid\phi,Z]$, $b_y=\Pr(Y=y\mid\phi,Z)$
(both conditional on $\phi=c$) for the two $K$-vectors on the right of
Eq.~\eqref{eq:sensitivity-exact}, and let
\begin{equation}
\rho=\frac{\sum_{y=1}^K\operatorname{Cov}(a_y,b_y\mid c)}
{\sqrt{\sum_{y=1}^K\operatorname{Var}(a_y\mid c)}\sqrt{\sum_{y=1}^K\operatorname{Var}(b_y\mid c)}}
\in[-1,1]
\end{equation}
be their aggregate correlation. Then
\begin{align}
\left|\Delta R_c(\phi)-\mathbb E[\Delta R_{(\phi,Z)}\mid\phi=c]\right|
&\ \le\ |\rho|\sqrt{\textstyle\sum_{y=1}^K\operatorname{Var}(a_y\mid c)}\,
\sqrt{\textstyle\sum_{y=1}^K\operatorname{Var}(b_y\mid c)}\notag\\
&\ \le\ |\rho|\sqrt{\textstyle\sum_{y=1}^K\operatorname{Var}(H_X(y)\mid\phi=c)}\,
\sqrt{\textstyle\sum_{y=1}^Kp_y(c)\bigl(1-p_y(c)\bigr)},
\label{eq:sensitivity-bound}
\end{align}
where $p_y(c)=\Pr(Y=y\mid\phi=c)$; the second inequality holds termwise
($\operatorname{Var}(a_y\mid c)\le\operatorname{Var}(H_X(y)\mid c)$ and
$\operatorname{Var}(b_y\mid c)\le p_y(c)(1-p_y(c))$ by the law of total variance, since
conditioning on more variables cannot increase the variance of a conditional
expectation) followed by monotonicity of the square root applied to each sum
separately -- a product of two square-root sums, not a sum of per-label square roots,
which the law of total variance alone would not license.
\end{proposition}

\begin{corollary}[Free bound]\label{cor:sensitivity-crude}
Each $a_y$ ranges over $[-M,M]$ as a conditional expectation of $H$, so
$\operatorname{Var}(a_y\mid c)\le M^2$ and $\sum_y\operatorname{Var}(a_y\mid c)\le KM^2$.
The label probabilities satisfy $\sum_yb_y=1$, which gives $\sum_yb_y^2\le1$ and, by
Cauchy--Schwarz, $\sum_y(\mathbb E b_y)^2\ge1/K$; hence
\begin{equation}
\sum_{y=1}^K\operatorname{Var}(b_y\mid c)
=\mathbb E\Bigl[\textstyle\sum_yb_y^2\Bigr]-\sum_y\bigl(\mathbb E b_y\bigr)^2
\ \le\ 1-\frac1K .
\label{eq:simplex-var}
\end{equation}
Substituting both into Eq.~\eqref{eq:sensitivity-bound} gives the data-free bound
\begin{equation}
\left|\Delta R_c(\phi)-\mathbb E[\Delta R_{(\phi,Z)}\mid\phi=c]\right|
\ \le\ |\rho|\,M\sqrt{K-1}.
\label{eq:sensitivity-crude}
\end{equation}
\end{corollary}

Earlier versions of this corollary bounded $\sum_y\operatorname{Var}(b_y\mid c)$ by
Popoviciu's inequality applied label by label, giving $K/4$ and hence $|\rho|KM/2$. That
discards Eq.~\eqref{eq:simplex-var}: the $b_y$ are not $K$ unconstrained quantities in
$[0,1]$ but a point on the simplex, and $K/4$ is unattainable for $K>2$. At VG150's
$K=50$ and $M=2$ the corrected free bound is $14.00|\rho|$ rather than $50.00|\rho|$, a
factor of $3.57$.

An analyst willing to bound $|\rho|\le\bar\rho$ -- how strongly an unrecorded variable
could jointly move a model's within-cell score gain and the within-cell label
distribution -- obtains a numeric limit on how far $\widehat{\Delta R}_c$ can be from
what a fully adjusted audit would find, without observing or modeling $Z$, in the spirit
of \citet{rosenbaum2002observational} and recent partial-correlation sensitivity models
for unmeasured confounding \citep{cinelli2020making}.
Because $\bar\rho$ is a correlation rather than a domain-specific odds or risk ratio, it
can be elicited the same way across benchmarks.

Both forms of the bound are evaluated on this paper's own Visual Genome estimates in
Appendix~\ref{app:sensitivity}. The short version is that the free bound is directionally
correct but far too loose to certify anything, while the sharp per-cell form is
\emph{less} reassuring than the free one: an unrecorded covariate whose aggregate
correlation with both channels reaches $0.061$ would account for the entire covariance gain
reported in Section~\ref{sec:results}. We therefore do not claim those gains are robust to
confounding, and Section~\ref{sec:sggaudit}'s audit is unaffected only because its estimate
is already indistinguishable from zero.

\begin{theorem}[Asymptotic normality of the covariance-gain estimator]\label{thm:clt}
Suppose (i) $H_x(y)\in[-M,M]$ for all $x,y$; (ii) examples are partitioned into clusters
$g=1,\dots,G$ (e.g.\ images) that are mutually independent, with cluster sizes having a
finite third moment and $\max_g|g|=o(\sqrt{G})$; (iii) the limiting variance
$\sigma^2=\lim_{G\to\infty}\operatorname{Var}\!\bigl(\sqrt{N_*}(\widehat{\Delta R}-\theta)\bigr)$
exists and is strictly positive, where $\theta$ is the identified-subpopulation estimand
of Appendix~\ref{app:influence}; and (iv) $\phi$ takes values in a fixed, finite set,
each with strictly positive limiting population mass, so that every eligible cell's size
$n_c$ diverges almost surely as $N\to\infty$. Then
\begin{equation}
\sqrt{N_*}\bigl(\widehat{\Delta R}-\theta\bigr)\ \xrightarrow{d}\ \mathcal N(0,\sigma^2),
\end{equation}
and the plug-in sandwich variance $\widehat\sigma^2=N_*\dfrac{G}{G-1}\sum_{g}\bigl(\sum_{i\in g}\widehat\psi_i/N_*\bigr)^2$
of Appendix~\ref{app:influence} is consistent for $\sigma^2$.
\end{theorem}

Appendix~\ref{app:clt} discusses what condition (iv) costs --- it forces every cell to be
eventually identified, so the theorem does not describe the finite-sample regime our own
applications sit in --- and states the additional condition on the number of cells that the
dataset-level remainder requires. What licenses reading the reported intervals at finite $N$
is the coverage simulation of Section~\ref{sec:simulation}, not this theorem.
